\cotefile{136}
\scanpage{127}
a) Chacun des $A^n(-)$ est muni d'une structure de $k$-coalgèbre \inserted{\unsure{(alternée)}} graduée codifférentielle (codifférentielle de degré $-1$), avec immersion $\overline S\hookrightarrow A^n(-)$ qui est compatible avec tout (\unsure{prend|prenant|provenant} sur $\overline S$ la gradation nulle, et la différentielle nulle). et \deleted{\unsure{une}}\inserted{un} syst. de copuissances divisées sur $\unsure{A^*|A^n}(\;)^+=A^n(\;)/\overline S$, compatibles avec graduation et différentielle.

b) Pour $n$ variable, \unsure{les} $A^n(\;)$ forment un objet $1/2$ \inserted{\uline{co}}simplicial (et $=$ simplicial) $A^{\unsure{*|\bullet}}(S,S^+,\gamma_*,t)$ dans la cat. des \inserted{objets munis \deleted{\unsure{de la}}de structure} (coAlg\unsure{....}) décrite dans a).

c) On réconstitue $A_*(S,S^+,\gamma_*,t)$ \inserted{$= k'\longmapsto A_*(\unsure{k|S},\unsure{\gamma_*},t)(k')$} avec \unsure{ttes} ses structures à l'aide des $A^*(S,S^+,\gamma_*,t)$, comme
\[
 A_*(S,S^+,\gamma_*,t)=V_k\bigl(A^*(S,S^+,\gamma_*,t)\bigr)
\]

d) Soit $X$ un ens. $1/2$ simplicial, considérons
\[
\begin{aligned}
 k'&\longmapsto \mathrm{Hom}\bigl(X,A_{\unsure{*|\bullet}}(-)(k')\bigr)
 \inserted{\deleted{\unsure{*}}\mathrel{\overset{\mathrm{dfn}}{=}}C_{\mathrm{DR}}^*(X;S^{k'},(S^{k'})^+,\gamma_*,t)},\\
 &\simeq\varprojlim_{\sigma\ \mathrm{ds}\ \Delta/X}
 \mathrm{Hom}\bigl(\sigma,A_{\bullet}(\;)(k')\bigr)\\
 &\phantom{\simeq}\hspace{1em}\underset{\substack{\text{cat. des simplexes}\\\text{types }(\Delta_i)_{i\geq0}\\\Delta_i=[0,i]}}{\downarrow}\\
 &\simeq\varprojlim_{\sigma\ \mathrm{ds}\ \Delta/X}
 \underbrace{A_{\deleted{\unsure{\mathrm{dim}}}\sigma}(\;)(k')}
 _{\mathrm{Hom}_k(A_{\deleted{\unsure{\mathrm{dim}\,\sigma}}}^{\mathrm{dim}\,\sigma}(-),k')}\\
 &\simeq\mathrm{\unsure{hom|Hom}}_k
 \bigl(\varinjlim_{\sigma\,\unsure{\in\mathrm{Ob}\ \mathrm{de}|\mathrm{ds}|\mathrm{dans}}\ {\Delta/X}}
 A_{\deleted{\unsure{\mathrm{dim}}}}^{\mathrm{dim}\,\sigma}(-),k'\bigr)\\
 &\simeq\mathrm{Hom}_k\bigl(C_*^{\mathrm{DR}}(X;S,S^+,\gamma_*,t),k'\bigr)
\end{aligned}
\]
% archive sheet number: 126

\scanpage{128}\mspage{- 118 -}\hfill IX

de type fini. On constate alors que l'homomorphisme (9.5.6) est l'homomorphisme canonique qui exprime l'action du sous-groupe d'inertie $I\subset\pi$ sur $X_\eta(\overline K)$, défini par la méthode de 9.2. Pour le vérifier, on se ramène immédiatement à la situation de 9.4, en y supposant que $n$ est égal à une puissance de $\ell$, et que $R$ est un $\Lambda$-Module localement constant (tout comme $Q$ et $P$), de sorte qu'il en est également de même de $E$, $F$ et de $X_\eta$. La vérification de l'identité des deux définitions est alors laissée au lecteur.

9.6.\quad Nous sommes maintenant en mesure de définir l'accouplement canonique annoncé (9.1.2), sans restriction sur $\ell$, de telle façon que pour $\ell\neq p$ cette définition \unsure{coïncide|coîncide} avec celle donnée dans 9.2. Pour ceci, on applique les considérations de 9.5 à la situation suivante :
\[
\left\{\begin{aligned}
 &P=T_\ell(A^{\unsure{o|\circ}})^t=\underline{\check M}'_\ell(1),\quad Q=M_\ell,\quad R=T_\ell(A^o)^f/T_\ell(A^o)^t\ ,\\
 &\hspace{2em}F=T_\ell(A^o)^f,\quad E=\unsure{,}D(F')=D(T_\ell(A'^{o})^f)\ ,
\end{aligned}\right.
\tag*{(9.6.1)}
\]
où dans la dernière formule $D$ désigne le dual de Cartier au sens des pro-$\ell$-groupes de Barsotti-Tate, i.e $D(Y)$ est le système \uline{projectif} des $D(Y(n))$, pour les morphismes de transition habituels. La structure de $F$ comme extension de $R$ par $P$ est claire. On a de même sur $F'=T_\ell(A'^{o})^f$ une structure d'extension de $R'$ par $P'=T_\ell(A'^{o})^t=\underline{\check M}_\ell(1)$. Appliquant la dualité de Cartier, on en conclut que $E=D(F')$ est muni d'une structure d'extension de $D(P')=M_\ell$ par $D(R')$. Or $D(R')$ est canoniquement isomorphe à $R$ en vertu de 5.5 (prouvé dans 7.4.6), ce qui précise la structure de $E$ comme extension de $Q$ par $R$.
% archive sheet number: 127
% stamp: numbered label 436; 436 also faintly visible above the lower strip.

\scanpage{129}
\unsure{ou|où}
\[
\boxed{C_*^{\mathrm{DR}}(X;S,S^+,\gamma_*,t)
 \mathrel{\overset{\mathrm{dfn}}{=}}\varinjlim_{\sigma\ \mathrm{ds}\ \Delta/X}
 A^{\mathrm{dim}\,\sigma}(S,S^+,\gamma_*,t)}
\]

Ici $C_*^{\mathrm{DR}}(X;S,S^+,\gamma_*,t)$ est une $k$-coalgèbre graduée (alternée) codiff. (codiff. de degré $-1$), \deleted{à copuissances divisées \unsure{sur}}\inserted{avec un} $\overline S\longrightarrow C_*^{\mathrm{DR}}(X,-)$ \inserted{\unsure{pr tte la structure}} (pas néc injectif, p. ex. si $X=\varnothing$ $C_*^{\mathrm{DR}}(X,-)=0\ \ldots$) et des puiss. divisées sur \deleted{\unsure{Cok}} $C_*^{\mathrm{DR}}(X,-)_+=\deleted{\unsure{\mathrm{Co}}}\unsure{\mathrm{CoKer}}(\overline S\longrightarrow C_*^{\mathrm{DR}}(X,-))$. Il est covariant en $X$ avec \unsure{ttes} ses structures, \inserted{\unsure{or|sur} $C_{\mathrm{DR}}^*(X,S^{k'},\ldots)$ est contra\unsure{variant|variantes} en id.}
% slanted
\margin{$X\longmapsto C_*^{\mathrm{DR}}(X,-)$ dépend de $X$ en commutant aux $\varinjlim$ quelconques (pas \unsure{seult} filtrantes)}

e) Si $(S,S^+,\gamma_*,t)$ varie sur $k$, \deleted{\unsure{il commute aux|varie avec $X$}} $(S^{k'},S^{+k'},\gamma_*^{k'},t^{k'})$ en dépend de façon contravariante, donc $C_{\mathrm{DR}}^*(X;S^{k'},S^{+k'},\gamma_*^{k'},t^{k'})$ aussi
\[
 \mathrm{Hom}(\unsure{\text{\textquotedblleft}}C_*^{\mathrm{DR}}(X,-),k')\ ,
\]
donc $C_*^{\mathrm{DR}}(X;S,S^+,\gamma_*,t)$ dépend de façon covariante en $(S,S^+,\gamma_*,t)$, ce qui est aussi évident sur l'écriture plus haut $C_*^{\mathrm{DR}}(X,-)=\varinjlim\ldots$, compte tenu que $A^*(S,S^+,\gamma_*,t)$ est covariant en $(S,S^+,\gamma_*,t)$. De plus, \deleted{il} comme foncteur en $(S,S^+,\gamma_*,t)$, il commute aux $\varinjlim$ \unsure{quelc.|filtrantes}

Commutation à l'extension de l'anneau de base.

f) Supposons que la coalgèbre $S$ soit graduée (structure de coalgèbre et graduation compatibles) que $\overline S$, soit \unsure{idem|itou}, $\overline S\to S$ compatible avec, les p.d. sur $S^+$ compatibles aux graduations et $t:S^+\to k$ \inserted{\unsure{hom.|homom.}} de degré $\delta$. Alors \deleted{\unsure{les}}\inserted{$\forall k'$}, $(S^{k'},S^{k'\,+},\gamma_*^{k'},t^{k'})$
% archive sheet number: 128

\scanpage{130}\mspage{- 117 -}\hfill IX

Pour tout entier $n$, l'homomorphisme correspondant $X(n)\longrightarrow Q(n)$, étant surjectif sur les fibres génériques, est surjectif ; comme $Q(n)$ est étale sur $n$, donc $X(n)\longrightarrow Q(n)$ est plat, il s'ensuit que c'est un épimorphisme. Donc $X\longrightarrow Q$ \uline{est un épimorphisme composant par composant}, Si donc $\overline F$ désigne le noyau de $X\longrightarrow Q$, il s'ensuit que $\overline F$ est un groupe de Barsotti-Tate et $X$ est une extension de $Q$ par $\overline F$. D'ailleurs l'isomorphisme $F_\eta\overset{\sim}{\longrightarrow}\overline F_\eta$ se prolonge, par le théorème de TATE, en un isomorphisme $F\simeq\overline F$. Par suite, $X$ apparaît comme une extension de $Q$ par $F$. Considérons alors le quotient $\overline E=X/P$, on voit encore, par application du théorème de TATE, qu'on a un isomorphisme $E\simeq\overline E$. Il s'ensuit que $X$ est une extension panachée de $E$ par $F$ prolongeant l'extension panachée $X_\eta$, cqfd.

\uline{Remarques} 9.5.6.4.\quad a)\quad Bien entendu, la dernière assertion de 9.5.6 devrait être valable sans restriction sur $K$, puisqu'il devrait en être ainsi du théorème de Tate invoqué.

b)\quad On peut donner une autre démonstration du critère de bonne réduction 5.10, en se ramenant comme dans 5.10 au cas où $A_K$ est à réduction semi-stable, et en appliquant alors 9.5.6 à $T_\ell(A_K)$ considérée comme extension panachée (cf. 9.6), qui implique $u_\ell=\unsure{0|O}$, et utilisant 11.6.2 a) plus bas qui implique que l'on a $T_\ell(A_K)^t=0$.

9.5.8.\quad Lorsque $\ell\neq p$, les groupes de Barsotti-Tate envisagés sur $S$ resp. sur $K$ s'interprètent simplement en termes de faisceaux $\ell$-adiques constants tordus sur $S$, ou de façon équivalente, en termes de représentations de $\pi_{\unsure{o|0}}=\mathrm{Gal}(\overline k/k)$ resp. de $\pi=\mathrm{Gal}(\overline K/K)$ sur des $\mathbb Z_\ell$-modules libres
% archive sheet number: 129
% stamp: numbered label 435.

\scanpage{131}
est une alg. \deleted{\unsure{à p. div.}} \uline{graduée}, avec idéal $(S^{k'})^+$ \uline{hom} à p. div. compatibles aux graduations, et $t^{k'}$ est un élt de degré $\delta$ du dit.

Alors les
\[
\begin{aligned}
 A_n(\;)(k')&=S^{k'}\{\{X_0,-,X_n\}\}[dX_0,-,dX_n]\\
 &\qquad\big/\bigl(\textstyle\sum X_i-t,\sum dX_i\bigr)_{\mathrm{pd}}
\end{aligned}
\]
\inserted{une alg.} est \uline{bigraduée}, par le degré extérieur déjà envisagé ($2^{\text{ème}}$ degré) \inserted{(nul sur $S^{k'}$)} et le degré total défini par
\[
\left\{\begin{aligned}
 &\deg\mathrm{\,tot}\,X_{r}=\deg\mathrm{\,tot}\,dX_r=\delta\quad\deleted{\unsure{\mathrm{degré}}}\\
 &\deg\mathrm{\,tot}\,X_r^{(\unsure{n|i})}=\delta,\quad
 \deg\mathrm{\,tot}_{A_n}\lambda=\deg\mathrm{\,tot}_{S^{k'}}\lambda\quad\big|\ \lambda\in S^{k'}
\end{aligned}\right.
\]
La diff. \deleted{\unsure{applica}} respecte le degré total, $A_n(\;)(k')$ est \unsure{pour} le degré total une \uline{alg.} graduée sur $S^{k'}$ graduée comme on a dit ; les puiss. divisées sont compatibles aux deux degrés.

\deleted{\unsure{Par con}}\inserted{Cette} structure supplémentaire (deg total) est fonctorielle en $k'$, donc passe à $A^*(S,S^+,\gamma_*,t)$. Elle passe par suite aux $C_*^{\mathrm{DR}}(X;S,S^+,\gamma^*,t)$, et on retrouve la structure déjà connue de $C_{\mathrm{DR}}^*(X;S^{k'},(S^{k'})^+,\gamma^{+k'},t^{k'})$ en prenant le $\mathrm{Hom}_k$ du \unsure{C} $C_*^{\mathrm{DR}}(X;S,S^+,\gamma^*;t)$ dans $k'$ \unsure{---|.}

g) On prendra maintenant, \deleted{\unsure{par|pour}}
\[
\begin{aligned}
 S&=\deleted{\unsure{\mathrm{df}}} S_k=k[u] & \text{donc }S^{k'}&=k'\{\{T\}\}\\
 \overline S&=\overline S_k=k & \text{donc }\overline S^{k'}&=k'\\
 &\inserted{\text{avec structure de coalg ordinaire, \deleted{\unsure{donn}}\unsure{donnant} sur }S^{k'}\text{ la structure connue}}&&\\
 \overline S&\longrightarrow S\quad\text{l'inclusion} &
 \text{donc }S^{k'}&\unsure{-|\longrightarrow}\overline S^{k'}\quad\text{l'augm. ordinaire}\\
 k&\unsure{-|\longrightarrow}k[u] & k'\{\{\unsure{\sigma|T}\}\}&\longrightarrow k'\\
 S^+&=k[u]/k\simeq k[u]^+ &&
\end{aligned}
\]
% archive sheet number: 130

\scanpage{132}\mspage{- 116 -}\hfill IX

\uline{dans la catégorie des extensions panachées de} $E_\eta$ \uline{par} $F_\eta$. \uline{Pour tout groupe de \unsure{Barsotti-Tate|Barootti-Tate} extension panachée} $X_\eta$ \uline{de} $E_\eta$ \uline{par} $F_\eta$, \uline{on considère l'accouplement correspondant}
\[
 Q_{\unsure{o|0}}\otimes Q_o^*\longrightarrow\mathbb Z_{\ell s}\ , \tag*{(9.5.6.2)}
\]
\uline{induit par l'accouplement du type} (9.5.5) \uline{défini sur le} hensélisé \uline{de} $S$. \uline{Ceci posé}, \uline{le foncteur} (9.5.6.1) \uline{est pleinement fidèle}, \uline{et son image essentielle est formée des groupes de \deleted{\illegible{1}}Barsotti-Tate} $X_\eta$ \uline{extensions panachées de} $E_\eta$ \uline{par} $F_\eta$ \uline{telles que l'accouplement} (9.5.6.2) \uline{correspondant soit nul}. \uline{De plus}, \uline{lorsque} $\mathrm{car}.K=0$ (\uline{de sorte qu'on peut appliquer le théorème de TATE} [33 th. 4]), \uline{alors} $X_\eta$ \uline{appartient à l'image essentielle si et seulement si} $X_\eta$ \uline{se prolonge en un groupe de Barsotti-Tate sur} $S$.

9.5.6.3.\quad Nous laissons au lecteur le soin d'étendre cet énoncé au cas où on suppose seulement que $S$ est un schéma noethérien régulier connexe de dimension 1.

La première assertion de 9.5.6 est une conséquence immédiate de 9.4.6, appliqué aux composants des groupes de Barsotti-Tate envisagés. Il reste à prouver que si un groupe de \deleted{\illegible{1}}Barsotti-Tate $X_\eta$, extension panachée de $E_\eta$ par $F_\eta$, se prolonge en un groupe de Barsotti-Tate $X_\eta$ sur $S$, il se prolonge aussi comme extension panachée. Or en vertu du théorème de TATE [33], l'homomorphisme $X_\eta\longrightarrow Q_\eta$ se prolonge de façon unique en un homomorphisme
\[
 X\longrightarrow Q\quad .
\]
% archive sheet number: 131
% stamp: numbered label 434.

\scanpage{133}
On prend syst. de copuiss. divisées sur $k[u]^+=S^+$ correspondant à celles \inserted{connues} \unsure{de|sur} $k'\{\{T\}\}^+$.

On prend $t\in k\unsure{[\{T\}\}|\{\{T\}\}|[\{T\}]}$ donné par $t=T$, i.e.
\[
 t(u^n)=\left\{\begin{aligned} &1\ \text{si }n=1\\ &0\ \text{si }n\neq1\end{aligned}\right.
\]
On a donc ici
\[
\begin{aligned}
 A_{\unsure{*|\bullet}}(S,S^+,\gamma^*,t)(k')&\simeq A_*(S^{k'},S^{+k'},\gamma_*^{k'},t^{k'})
 \simeq A_*(k'\{\{T\}\},k'\{\{T\}\}^+,\gamma^{k'},T)\\
 &\hspace{8em}\simeq A_*(k')^\wedge
\end{aligned}
\]
on posera simplement
\[
 \deleted{\unsure{A_k^*}}\, A_{k\deleted{\unsure{*}}}^*
 = A_{\deleted{\unsure{*}}}^*(S_k,\overline S_k,\gamma_k^*,t_k)
\]
de sorte que
\[
 A_*(k')^\wedge\simeq\mathrm{Hom}_k(A_k^*,k')
\]
``avec toutes les structures''. Notons qu'ici on est dans le cas où $S$ \deleted{\unsure{et}}, $\overline S$ et $\overline S\unsure{\to|-}S$ sont gradués (graduation ordinaire de $k[u]$, correspondant à la graduation ordinaire sur les $k'\{\{T\}\}$) et $t=T$ de \inserted{hom} degré $1$, donc \unsure{on|on a} \deleted{\unsure{A(k)|\text{on a}}} $A_k^*$ admet une bigraduation. Celle ci passe aux $\deleted{\unsure{C_*^{\mathrm{DR}}(|A_k^*}}$
\[
\begin{aligned}
 &\deleted{\unsure{C_*^{\mathrm{DR}}}}(X;k)\\
 &C_*^{\mathrm{DR\,pd}}(X,k)\mathrel{\overset{\mathrm{dfn}}{\simeq}}
 C_*^{\mathrm{DR}}(X,k[u];k[u]^+,\gamma_k^*,t_k)
\end{aligned}
\]
et redonne \inserted{(avec les autres)} par $\mathrm{Hom}_k(C_*^{\mathrm{DR\,pd}}(X,k),k')$ les structures connues (y compris bigraduation) sur le complexe de cochaines de De Rham $C_{\mathrm{DR\,pd}}^*(X,k')$.

En fait tous les $C_*^{\mathrm{DR\,pd}}(X,k)$ ($k$ variable)
% archive sheet number: 132

\scanpage{134}\mspage{- 115 -}\hfill IX

enfin donnée sur $\eta$ un \unsure{pro}-$\ell$-groupe de Barsotti-Tate $X_\eta$ qui soit \uline{extension panachée} de $E_\eta$ par $F_\eta$, par quoi nous entendons qu'il est muni d'une filtration en trois crans (en tant que pro-$\ell$-groupe de Barsotti-Tate), donnant lieu à des extensions partielles munies d'isomorphismes avec $E_\eta$ et $F_\eta$ respectivement. Il s'ensuit que pour tout $n$, $X_\eta(n)$ est muni d'une structure d'extension panachée de $E(n)_\eta$ par $F(n)_\eta$. Nous pouvons lui faire correspondre alors un accouplement du type (9.4.1) :
\[
 c(X(n)_\eta):Q(n)\otimes Q^*(n)\longrightarrow\Lambda_n=\mathbb Z/\ell^{n+1}\mathbb Z\ ,\tag*{(9.5.4)}
\]
où les $Q^*(n)$ sont les groupes finis étales sur $S$ définis par
\[
 P(n)\simeq D(Q^*(n))\ ,\quad\text{i.e.}\quad Q^*(n)=D(P(n))\ ,
\]
$D$ désignant encore la dualité de Cartier. Les $Q^*(n)$ forment un pro-$\ell$-groupe de façon évidente, soit $Q^*$. Nous laissons au lecteur le soin de vérifier que pour $n$ variable, les homomorphismes (9.5.4) définissent un homomorphisme de systèmes projectifs, i.e. un homomorphisme de faisceaux $\ell$-adiques constants tordus sur $S$ :
\[
 c(X_\eta):Q\otimes Q^*\longrightarrow\mathbb Z_{\ell S}\quad . \tag*{(9.5.5)}
\]

Signalons en passant le résultat suivant (qui ne sera pas utilisé dans la suite) :

\uline{Proposition} 9.5.6. \uline{Les données} $P,Q,R,E,F$ \uline{sont les mêmes que ci-dessus}, \uline{sauf qu'il est inutile de supposer le trait} $S$ \uline{hensélien.} \uline{On considère le foncteur}
\[
 X\longmapsto X_\eta \tag*{(9.5.6.1)}
\]
\uline{de la catégorie des extensions panachées de Barsotti-Tate de} $E$ \uline{par} $F$
% archive sheet number: 133
% stamp: numbered label 433.

\scanpage{135}
se déduisent des
\[
 C_*^{\mathrm{DR\,pd}}(X)\mathrel{\overset{\mathrm{dfn}}{=}}C_*^{\mathrm{DR\,pd}}(X;\mathbb Z)
\]
par la formule
\[
 C_*^{\mathrm{DR\,pd}}(X,k)\simeq C_*^{\mathrm{DR\,pd}}(X)\otimes_{\mathbb Z}k\ ,
\]
et \unsure{on a|d'où} de $=$
\[
 C_{\mathrm{DR\,pd}}^*(X,k)\simeq\mathrm{Hom}\bigl(C_*^{\mathrm{DR\,pd}}(X,k)\bigr)
\]

On définit de $=$, pour tt groupe abélien $M$ et \unsure{fonctoriellmt|fonctoriellement} en $M$
\[
\begin{aligned}
 C_*^{\mathrm{DR\,pd}}(X;M)&\mathrel{\overset{\mathrm{dfn}}{=}}C_*^{\mathrm{DR\,pd}}(X)\otimes_{\mathbb Z} M\\
 C_{\mathrm{DR\,pd}}^*(X;M)&\mathrel{\overset{\mathrm{dfn}}{=}}\mathrm{Hom}\bigl(C_*^{\mathrm{DR\,pd}}(X),M\bigr)
\end{aligned}
\]
et les structures sur $C_*^{\mathrm{DR\,pd}}(X)$ impliquent une foule de structures multiplicatives pour \unsure{ce foncteur|ces foncteurs} en $M$ \unsure{\ldots}.
% archive sheet number: 134

\scanpage{136}\mspage{- 114 -}\hfill IX

qui est injective (puisque la suite exacte $0\longrightarrow V^*\longrightarrow K^*\longrightarrow\mathbb Z\longrightarrow0$ est scindée), donc elle est injective, ce qui implique (grâce au \unsure{sorite} 9.3) que si $X_\eta$ et $Y_\eta$ sont isomorphes, il en est de même de $X$ et de $Y$. On gagne donc par a).

Evidemment a) et b) impliquent que le foncteur $X\longmapsto X_\eta$ est \uline{pleinement} fidèle.

c)\quad Soit $X_\eta$ une extension panachée sur $\eta$, prouvons qu'elle se prolonge à $S$ si et seulement si les accouplements (9.4.6.2) sont nuls. Ici encore, l'unicité déjà prouvée montre que la question est locale pour la topologie étale, ce qui nous\unsure{,} ramène encore au cas $S$ strictement local. Alors la conclusion est vraie par construction, comme on avait signalé juste avant 9.4.6.

Ainsi la démonstration de 9.4.7 est achevée.

9.5.\quad Supposons maintenant fixé un nombre premier $\ell$, des pro-$\ell$-groupes de Barsotti-Tate sur $S$,
\[
 P=(P(n))_{n\geq0}\ ,\ R=(R(n))_{n\geq0}\ ,\quad Q=(Q(n))_{n\geq0}\ ,\tag*{(9.5.1)}
\]
et des extensions de pro-$\ell$-groupes de Barsotti-Tate
\[
 0\longrightarrow P\longrightarrow F\longrightarrow R\longrightarrow0\ ,\quad
 0\longrightarrow R\longrightarrow E\longrightarrow Q\longrightarrow0\ ,\tag*{(9.5.2)}
\]
définissant donc pour tout entier $n\geq0$ des extensions de $\Lambda_n$-Modules où $\Lambda_n=\mathbb Z/\ell^{n+1}\mathbb Z$ :
\[
 0\longrightarrow P(n)\longrightarrow F(n)\longrightarrow R(n)\longrightarrow0,\quad
 0\longrightarrow R(n)\longrightarrow E(n)\longrightarrow Q(n)\longrightarrow\unsure{0|C}\ .\tag*{(9.5.3)}
\]
Nous supposerons $Q$ de type étale et $P$ de type multiplicatif. Supposons
% archive sheet number: 135
% stamp: numbered label 432; the same number appears in pencil above the strip.

\scanpage{137}
\deleted{ce qui est évident, car}
\[
 \deleted{\tau^*(F)\simeq|\alpha|^*(\sigma^*(F_{\unsure{i|,}}))\quad\ldots}
\]
\deleted{Comme les $\sigma^*$ commutent aux $\varprojlim$ finies et aux $\varinjlim$ \unsure{quelc.|quelconques}}

\diagram{\begin{tikzcd}[ampersand replacement=\&,column sep=small]
 \deleted{X'=\varprojlim X_{\mathcal U}}\arrow[rr]\arrow[dr] \& \& \deleted{X}\\
 \&\deleted{X_{\mathcal U}}\arrow[dr]\&\\
 \&\&\deleted{X}
\end{tikzcd}}
\[
 \deleted{\varinjlim\,\unsure{\mathcal F|F}/X_i}
\]
\diagram{\begin{tikzcd}[ampersand replacement=\&]
 \deleted{\unsure{Q_i}}\& \deleted{\Delta/X_i}\&\deleted{\Delta/X_j}\\
 \deleted{\unsure{q_i}}\quad\deleted{\unsure{\Omega}}\&\deleted{X_i}\arrow[r,hook]\&\deleted{X_j}\\
 \&\deleted{\unsure{F_i|\mathcal F_i}}\&\deleted{F_j}\arrow[l]
\end{tikzcd}}

\uline{Dualisation des hom de Stokes}

\diagram{\begin{tikzcd}[ampersand replacement=\&,column sep=small,row sep=small]
 {\pi_n:A_n^n(S,S^+,\gamma_{\deleted{\unsure{*}}}^*,t)(\unsure{k'|k})=A_n^\wedge(S^{k'},(S^{k'})^+,\gamma_*^{k'},t^{k'})^n}\arrow[d,equal]\\
 {V\bigl(S[Y_0,-,Y_n][\deleted{\unsure{Y_0}}dY_0,\unsure{-|{,}|\ldots},\deleted{\unsure{Y_n}}dY_n]^n\bigr)(k')}\arrow[d,"{\pi_{\deleted{\unsure{n}}}^n\quad S\ \text{\uline{hom}}}"]\\
 {\underline{V(S)(k')}=S^{k'}}
\end{tikzcd}}
\[
 (A_{\deleted{\unsure{*}}}^*)_n\qquad
 S\xrightarrow{\pi_n} S[Y_0,-,Y_n][\unsure{dY_0|Y_0|y_0},-,Y_n]^{(\unsure{n,*|n})}
\]

1) C'est l'identité si $n=0$

2) Stokes : comm. dans
\margin{$\unsure{K|\pi}_{\deleted{\unsure{n}}}^n(d\omega)=K_{\deleted{\unsure{n}}}^{n-1}(\partial\omega)$\\
$\Updownarrow\qquad\downarrow$\\
$\underbrace{d_{A^{\unsure{*}}}}_{\text{degré ext. }-1}\pi_n=\underbrace{\partial_{\mathrm{ss}}}_{\text{deg ss }+1}\pi_{n-1}$}
\margin{\deleted{$\unsure{K_n(d}$\\
$\underbrace{d_{A^*}}_{\text{degré }\unsure{1|+1}}\!K_n=\unsure{\partial}K_{n+1}\underbrace{\partial_{\mathrm{ss}}}_{\text{deg }+1}$}}
\diagram{\begin{tikzcd}[ampersand replacement=\&]
 S\arrow[r,"{\unsure{k_n|\pi_n|\kappa_n}}"]\arrow[dr,phantom,"\cdot" description]\arrow[d,"{\pi_{n-1}}"']\&{(A^{\unsure{n|*}})_n}\arrow[d,"{d_{A^*}}"]\\
 {(A^{\unsure{n-1|*}})_{n-1}}\arrow[r,"{\partial_{\mathrm{ss}}}"']\&{(A^n)_{n-1}}
\end{tikzcd}}
$S\deleted{\unsure{k'|\otimes k'}}$
\[
\begin{aligned}
 C_*^{\mathrm{sing}}(X,\unsure{S|s})&\longrightarrow C_*^{\mathrm{DR}\deleted{\unsure{p}}}(X;S,S^+,\gamma_*,t)\\
 C_*^{\mathrm{sing}}(X,k[u])&\longrightarrow C_*^{\mathrm{DR}\unsure{\mathrm{pd}}}(X,k)
\end{aligned}
\]
% archive sheet number: 136

\scanpage{138}\mspage{- 137 -}\hfill IX

sur $K$, qui jouera dans le cas présent le même rôle que 11.2 dans le cas traité précédemment :

\uline{Théorème} 11.5. \uline{Supposons} $S$ \uline{hensélien}, \uline{et que} $A_K$ \uline{ait réduction semi-stable sur} $S$. \uline{Soit} $\ell$ \uline{un nombre premier}, \uline{et considérons l'accouplement de monodromie} (9.1.2), \uline{d'où un homomorphisme canonique}
\[
 \widetilde u_\ell:\underline M_\ell\longrightarrow\underline{\check M}'_\ell\quad .
\]
\uline{Soit} $\Phi(\ell)=\mathrm{coker}\,\widetilde u_\ell$ ; \uline{alors on a un isomorphisme canonique}
\[
 \Phi_{\unsure{o|0}}(\ell)\simeq\Phi(\ell)\times_S s\quad .\tag*{(11.5.1)}
\]

De cet énoncé résulte immédiatement la définition d'un accouplement parfait (11.4.1), puisque on trouve, en appliquant 11.5 à $A'_K$ au lieu de $A_K$ et utilisant (10.2.2), qu'on a aussi un isomorphisme canonique
\[
 \Phi'_o(\ell)\simeq\Phi'(\ell)\times_S s\quad,\quad\text{où }\Phi'(\ell)=\mathrm{coker}({}^t\widetilde u_\ell:\underline M'_\ell\longrightarrow\underline{\check M}_\ell)\quad .
\]

\uline{Remarques} 11.5.2. a) Le théorème 11.5 implique notamment que $\widetilde u_\ell$ est une isogénie i.e. que $u_\ell$ est séparant, y compris dans le cas $\ell=p$, où cette assertion n'est pas triviale sur la définition de 9.6. On notera que nous n'\deleted{\illegible{1}}utilisons pas ici le théorème (nettement plus profond) 10.4, qui sera prouvé au par. suivant.

b)\quad Utilisant 10.4, qui nous fournit un accouplement de monodromie $u:\underline M\otimes\underline M'\longrightarrow\mathbb Z_S$, et désignant par $\Phi$ le $\underleftarrow{\text{\deleted{\unsure{o|c|g}}\handwritten{\inserted{\unsure{c}}}onoyau}}$ de l'homomorphisme correspondant
\[
 \widetilde u:\underline M\longrightarrow\underline{\check M}'\quad ,
\]
\margin{\unsure{2}}
% archive sheet number: 137
% stamp: numbered label 455.

\scanpage{139}
\diagram{\begin{tikzcd}[ampersand replacement=\&,column sep=small]
 {}\&{}\arrow[l,two heads,tail]\arrow[r,tail]\&{}
\end{tikzcd}}
\diagram{\begin{tikzcd}[ampersand replacement=\&,column sep=small]
 {A_n^{\overset{\nearrow\,\text{deg ext}}{0}}\simeq}
 \& {k\{\{T\}\}\{\{X_0,-,X_n\}\}/(\sum X_i-T)_{\mathrm{pd}}}\arrow[d,equal]
 \&{k\{\{T\}\}\{\{X_0,--,X_n\}\}}\arrow[l,two heads,tail,"\text{épi}"']\arrow[d,equal]\\
 \& {k\{\{X_0---,X_n\}\}}
 \&{k\{\{T,X_0,---,X_n\}\}}\arrow[l,two heads,tail,"{\substack{X_i\,\leftarrow\!\shortmid\, X_i\\ \sum_{0\leq i\leq n} X_i\,\leftarrow\!\shortmid\,T}}"]
\end{tikzcd}}
$\deleted{\unsure{\twoheadleftarrow}}$

duals
\diagram{\begin{tikzcd}[ampersand replacement=\&,column sep=small]
 {(A^n)_{\underset{\downarrow\,\text{deg ext}}{0}}\simeq}
 \& {k[u+Y_0,u+Y_1,-,u+Y_n]}\arrow[r,tail,"\mathrm{mono}"]\arrow[d,equal]
 \& {k[u][Y_0,-,Y_n]}\arrow[d,equal]\\
 \& {k[Z_0,Z_1,--,Z_n]}\arrow[r,tail,"{Z_i\longmapsto u+Y_i}"']
 \& {k[u,Y_0,--,Y_n]}
\end{tikzcd}}
\margin{[l'image est l'intersection des noyaux des op. diff
$\bigl(\frac{\partial}{\partial Y_0}+\,\unsure{--|\ldots+}\,\frac{\partial}{\partial Y_n}-\frac{\partial}{\partial u}\bigr)^{(i)}$]}
\[
\begin{aligned}
 \langle X_0^{(i_0)}\ldots X_n^{(i_n)},Z_0^{j_0}\ldots Z_n^{j_n}\rangle
 &=\prod_{0\leq\alpha\leq n}\delta_{i_\alpha,j_\alpha}\\
 &=\left\{\begin{aligned}&0\ \text{si }(i_\alpha)\neq(j_\alpha)\\&1\ \text{si }(i_\alpha)=(j_\alpha)\end{aligned}\right.
\end{aligned}
\]
\[
\begin{aligned}
 \langle T^{(\lambda)}X_0^{(i_0)}\ldots X_n^{(i_n)},u^\mu Y_0^{j_0}\ldots Y_n^{j_n}\rangle
 &=\delta_{\lambda,\mu}\delta_{\unsure{(i_\alpha)(j_\alpha)|(i_\alpha),(j_\alpha)|i_\alpha,j_\alpha}}\\
 &=\left\{\begin{aligned}&0\ \text{si }(\lambda,(i_\alpha))\neq(\mu,(j_\alpha))\\&1\ \text{si }\ -\ =\ -\end{aligned}\right.
\end{aligned}
\]

\diagram{\begin{tikzcd}[ampersand replacement=\&,column sep=small]
 {A_n^{*=*}\simeq}\&{\bigwedge_k^*[dX_0,--,dX_n]/(\sum dX_i)}
 \&{\bigwedge_k^*[dX_0,-,dX_n]}\arrow[l,two heads,tail,"\text{épi}"']
\end{tikzcd}}
\inserted{degré tot $=$ deg ext i.e. deg compl. $=0$}

duals
\diagram{\begin{tikzcd}[ampersand replacement=\&,column sep=small]
 {(A^n)_{*=*}\simeq}\&{\bigl(\bigwedge_k^*[y_0,--,y_n]\bigr)^{!}}\arrow[r,tail]\arrow[d,equal]
 \&{\bigwedge_k^*[y_0,---,y_n]}\\
 \&{\bigwedge_k^*(\underbrace{[y_0,-,y_n]^{!}}_{\substack{\text{ss-esp. de}\\ky_0\oplus--\oplus ky_n\\\text{formé des vect. dont la}\\ \text{somme des coord.}=0}})}\&
\end{tikzcd}}
\margin{l'image est le noyau de $i_X$, où $X$ est la forme sur $ky_0\oplus--\oplus ky_n$ ``somme des coordonnées'' et où $i_X$ est le \unsure{prod. int|produit int}}

\uline{NB} On a un isom can
\[
 \bigwedge_k^{\deleted{\unsure{n+1}}\,\unsure{n-1|n+1|n}}([y_0-y_n]')\simeq\bigwedge_k^n([y_0,-,y_n])\simeq k
\]
déduit de la suite exacte
\[
 \unsure{0|\cdot}\longrightarrow[y_0,-,y_n]'\longrightarrow[y_0,-,y_n]\longrightarrow k\longrightarrow0
\]
Donc
\[
 (A^n)_{\underset{\searrow}{n}}\mathrel{\overset{\text{iso can}}{\simeq}}(A^n)_{\underset{\swarrow}{0}}\simeq k[Z_0,-,Z_n]
 \qquad\text{\unsure{deg ext|deg cot}}
\]
donc
\[
 \pi_n:k[u]\longrightarrow k[Z_0,---,Z_n]
\]
% archive sheet number: 138

\scanpage{140}\mspage{- 135 -}\hfill IX

qui montre que $\Phi$ \uline{est canoniquement isomorphe au sous-groupe de torsion de} $H^1(I,U)$.

D'autre part, de la suite exacte (11.3.3) on déduit un isomorphisme, pour tout $n=\ell^r$ :
\[
 H^1(I,T_n)\simeq\mathrm{Hom}(I,T_n)\overset{\sim}{\longleftarrow}\mathrm{Hom}(T,T_n)\overset{\sim}{\longleftarrow}\mathbb Z/n\mathbb Z \tag*{(11.3.9)}
\]
(NB on fait opérer $I$ trivialement sur $T$, donc sur $T_n=T/nT$) ; d'où, pour tout $I$-$\mathbb Z_\ell$-Module fini $M$, si
\[
 M'=\mathrm{Hom}(M,T\otimes D_\ell) \tag*{(11.3.10)}
\]
désigne son ``dual de Cartier'', un accouplement canonique
\[
 H^i(I,M)\otimes H^{1-i}(I,M')\longrightarrow H^1(I,T\otimes D_\ell)\simeq D_\ell\ ,\tag*{(11.3.11)}
\]
où le dernier isomorphisme est déduit des isomorphismes (11.3.9) par passage à la limite inductive sur les puissances $n$ de $\ell$. Il est bien connu que ces accouplements sont des \uline{dualités parfaites} (cf. par exemple SGA 5 I 5). Prenant dans ces accouplements $M$ de la forme $U_n$, donc $M'=U'_n$, on trouve par passage à la limite un accouplement parfait
\[
 H^{\unsure{o|0}}(I,U\otimes D_\ell)\times H^1(I,U')\longrightarrow D_\ell\ ,\tag*{(11.3.12)}
\]
où le deuxième facteur du premier membre est un $\mathbb Z_\ell$-module de type fini. Un nouveau passage à la limite dans les accouplements (11.3.12) (pour des morphismes de transition de la forme $\ell^r.\mathrm{id}$) nous donne un deuxième accouplement parfait
\[
 H^o(I,U\otimes\mathbb Q_\ell)\times H^1(I,U'\otimes\mathbb Q_\ell)\longrightarrow D_\ell \tag*{(11.3.13)}
\]
de vectoriels de dimension finie sur $\mathbb Q_\ell$. Désignant par $i'$, $j'$ les
% archive sheet number: 139
% stamp: numbered label 453. Three broken vertical impressions recur at the left, without an identifiable text line.

\scanpage{141}
$S,S^+,\gamma^*,t$

\[
 (A_n^*)^n\simeq S\{\deleted{\unsure{\{}}X_0,-,X_n\deleted{\unsure{\}}}\}/(\textstyle\sum X_i-t)_{\mathrm{pd}}\otimes_S
 \underbrace{S[dX_0,-,dX_n]/(\textstyle\sum dX_i)}_{
 \substack{\text{en degré }n\\\text{c'est }\simeq S\\\text{avec la base}\\dX_0\wedge--\wedge dX_{n-1}}}
\]
\diagram{\begin{tikzcd}[ampersand replacement=\&]
 {\unsure{\simeq|\nearrow}}\&{}\\
 {S\{X_0,-,X_{n-1}\}[dX_0-dX_{n-1}]}\arrow[d,"\pi^n\ (S\text{-linéaire})"]\\
 S
\end{tikzcd}}
\[
 \pi^n(X_0^{(i_0)}--X_{n-1}^{(i_{n-1})}dX_0^{\deleted{\unsure{(i_0)}}}\wedge--\wedge dX_{n-1})
 =\deleted{\unsure{t^{(\sum i_\alpha)}\gamma_{i_\alpha}}}\,t^{\overbrace{(n+\sum i_\alpha)}^{\sum(i_\alpha+1)}}\gamma_{i_0,-,i_{n-1}}
\]
où $\gamma_{i_0--i_{n-1}}$ est un entier canonique
\[
 \gamma_{i_0--i_{n-1}}=\frac{(n+\sum\deleted{\unsure{l|d}}i_\alpha)!}{\prod_\alpha(\deleted{\unsure{l|d}}i_\alpha!)}
 \underbrace{\underset{\substack{\cdots\\ n}}{\int}\!\int_{\substack{\sigma_n=\{X_0,--,X_{n-\unsure{1}}\geq0\\X_0+--+X_{n-1}\leq1}}
 X_0^{i_0}--X_{n-1}^{i_{n-1}}\,dX_0\wedge--\wedge dX_{n-1}}_{0\leq\qquad\leq\frac1{n!}\gamma_{0,0,--,0}\mathrel{\overset{?}{=}}\frac1{n!}}
\]

On peut aussi calculer par récurrence sur $n$, en utilisant Stokes
\[
 \begin{aligned}
 &\pi^{\unsure{n|-}}\bigl(\underbrace{X_0^{(i_0)}--X_{n-1}^{(i_{n-1})}dX_0\wedge--\wedge dX_{n-1}}_{
 d\left(X_{n-1}^{(i_{n-1}+1)}dX_0^{(i_0+1)}-\deleted{\unsure{d}}-dX_{n-2}^{(i_{n-2}+1)}\right)}\bigr)\\
 &=\pi^{n-1}\bigl(\underbrace{\partial_{\mathrm{ss}}X_{n-1}^{(i_{n-1}+1)}d(X_0^{(i_0+1)})\,\ldots}_{
 \substack{X_{n-1}^{(i_{n-1}+1)}dY_0^{(i_0+1)}--\\d(X_{n-2}^{(i_{n-2}+1)})}}\bigr)
 \end{aligned}
\]
i.e. $\mathrel{/\!/}$ $\deleted{\unsure{\partial}}$
\[
 \begin{aligned}
 &\deleted{\unsure{\gamma_{i_0--i_{n-1}}=}}\pi_{n-1}\bigl(
 \underbrace{(t-(X_0+\cdots+X_{n-2}))^{(i_{n-1}+1)}}_{
 \substack{\displaystyle\sum_{\substack{\deleted{\unsure{l_0+--+l_{n-1}}}\\\underline\lambda+\underline{l_0}+--+\underline{l_{n-2}}=i_{n-1}+1}}
 (-1)^{\sum l_\alpha\deleted{\unsure{+d}}} X_0^{(l_0)}--X_{n-2}^{(l_{\unsure{n-1|n-2}})}t^{(\lambda)}}}
 X_0^{(i_0)}--X_{n-2}^{(i_{n-2})}dX_0\wedge--\wedge dX_{n-2}\bigr)\\
 &=\deleted{\unsure{(-1)^{i_{n-1}+1}}}\bigl(\pi_{n-1}\bigl(
 \sum_{\lambda+\sum_0^{n-2}l_\alpha=i_{n-1}+1}
 (-1)^{\sum l_\alpha\deleted{\unsure{+i_{n-1}+1}}}
 C_{i_0,l_0}\,C_{i_1,l_1}--C_{i_{n-2},l_{n-2}}\,
 \overset{\displaystyle C_{\lambda,\sum_0^{n-2}(i_\alpha+l_\alpha+1)}}{t^{(\lambda)}}
 X_0^{(i_0+l_0)}--X_{n-2}^{(i_{n-2}+l_{n-2})}dX_0\wedge--\wedge dX_{n-2}\bigr)\\
 &=t^{\deleted{\unsure{\sum}}(n+\sum i_\alpha)}
 \bigl[(-1)^{i_{n-1}+1}\deleted{\unsure{C|t}}
 \underbrace{\sum_{\underline\lambda+\sum\underline{l_\alpha}=i_{n-1}+1}(-1)^{\sum_\alpha l_\alpha}
 C_{i_0,l_0}--C_{i_{n-2},l_{n-2}}\,
 \gamma_{i_0+l_0,--,i_{n-2}+l_{n-2}}}_{\swarrow}
 \end{aligned}
\]
Donc
\[
 \begin{aligned}
 \gamma_{i_0--,i_{n-1}}={}&\deleted{(-1)^{i_{n-1}+1}}
 \sum_{\underline\lambda+\sum_0^{n-2}\underline{l_\alpha}=i_{n-1}+1}
 (-1)^{\sum_{\unsure{\alpha|0}}^{n-2}l_\alpha}
 C_{i_0l_0}--\deleted{\unsure{\}}}C_{i_{n-2},l_{n-2}}\\[-2pt]
 &\hspace{4em} C_{\lambda,\deleted{\unsure{n}}n-1+\sum_0^{n-2}(l_\alpha+i_\alpha)}\\
 &\hspace{1em}\gamma_{i_0+l_0,i_1+l_1,--,i_{n-2}+l_{n-2}}
 \end{aligned}
\]
% archive sheet number: 140
